
\subsection{弗洛尔理论中的连接轨道情形}
\label{subsec:8-1-3}

设 $(L^{(0)},L^{(1)})$ 是一对横截相交的相对 spin 拉格朗日
子流形（见定义~\ref{def:8-1-2}）。$L^{(0)},L^{(1)}$ 的弗洛尔
复形中的连接轨道，是一个伪全纯映射
\[
 u\colon\mathbb R\times[0,1]\longrightarrow M,
\]
它把 $\mathbb R\times\{i\}$ 映到 $L^{(i)}$，并满足
\[
 u(-\infty,*)=p^-,\qquad u(+\infty,*)=p^+.
\]
本章采用与第 2.3 节不同的记号：我们写
$\mathcal M^{\mathrm{reg}}(p^+,p^-)$，而不写
$\mathcal M^{\mathrm{reg}}(p^-,p^+)$。换言之，对
$u\in\mathcal M^{\mathrm{reg}}(p^+,p^-)$，有
$u(-\infty,*)=p^-$、$u(+\infty,*)=p^+$。这是因为，在我们的约定下，
前一种记号更适合定向问题，并且与第 8.4 节为定义滤过 $A_\infty$
代数而讨论的模空间定向相容。

考虑切丛 $TM$ 在投影 $M\times[0,1]\to M$ 下的拉回。
对 $p\in L^{(0)}\cap L^{(1)}$，选取 $T_pM$ 中定向拉格朗日
子空间的一条道路
\[
 \lambda_p\colon[0,1]\longrightarrow\Lambda^{\mathrm{ori}}(T_pM),
 \qquad
 \lambda_p(0)=T_pL^{(0)},\quad
 \lambda_p(1)=T_pL^{(1)},
\]
并选取丛
\[
 \widetilde\lambda_p
 =\bigsqcup_{t\in[0,1]}\{t\}\times\lambda_p(t)
 \longrightarrow[0,1]
\]
相对于端点处固定平凡化的一个平凡化；见第 2.2.2 小节。
于是，上述 $TM$ 拉回丛在
\[
\begin{aligned}
 K(L^{(0)},L^{(1)})
 :={}&(L^{(0)}\times\{0\})
 \cup\bigl((L^{(0)}\cap L^{(1)})\times[0,1]\bigr)\\
 &\cup(L^{(1)}\times\{1\})
\end{aligned}
\]
上有一个拉格朗日子丛
\[
 \mathcal L\longrightarrow K(L^{(0)},L^{(1)}).
\]

取 $st\in H^2(M;\mathbb Z_2)$，使得
\[
 w_2(TL^{(i)})=st|_{L^{(i)}}\qquad(i=0,1).
\]
如定义~\ref{def:8-1-2}，在 $M$ 的 $3$-骨架上选取满足
$w_2(V)=st$ 的定向向量丛 $V$。记
$P_{\mathrm{spin}}(L^{(i)},V)$ 为
$(TL^{(i)}\oplus V)|_{L^{(i)}_{[2]}}$ 的 spin 主丛；它逐纤维
双覆盖该丛的定向标架丛 $P_{SO}(L^{(i)},V)$。不妨假设
$L^{(0)}\cap L^{(1)}$ 包含在 $M$ 的 $3$-骨架中。

为了指定
\[
 (\mathcal L\oplus\operatorname{pr}^*V)
 \big|_{K(L^{(0)},L^{(1)})_{[2]}}
\]
上的 spin 结构，还需要附加资料；这里
$K(L^{(0)},L^{(1)})_{[2]}$ 是 $K(L^{(0)},L^{(1)})$ 的
$2$-骨架，而 $\operatorname{pr}$ 是投影 $M\times[0,1]\to M$。
设
\[
 \sigma\colon[0,1]\times\mathbb R^n
 \xrightarrow{\ \cong\ }\widetilde\lambda_p
\]
是一个平凡化。由于 $\lambda_p(i)=T_pL^{(i)}$，限制
$\sigma|_{t=i}$ 给出 $T_pL^{(i)}$ 的一个标架，从而得到嵌入
\[
 \sigma_{t=i,*}\colon SO(n)\times SO(V_p)
 \longrightarrow P_{SO}(L^{(i)},V)|_p
 \qquad(i=0,1).
\]
再取覆盖 $\sigma_{t=i,*}$ 的映射
\[
 \iota_i\colon
 \frac{\operatorname{Spin}(n)\times\operatorname{Spin}(V_p)}
      {\{\pm1\}}
 \longrightarrow P_{\mathrm{spin}}(L^{(i)},V)|_p,
\]
其中 $\{\pm1\}$ 对
$\operatorname{Spin}(n)\times\operatorname{Spin}(V_p)$ 对角作用。
现在考虑三元组 $(\sigma,\iota_0,\iota_1)$。

对平凡化 $\sigma$，选取覆盖 $\sigma_*$ 的丛同构
\[
 \widetilde\sigma\colon
 [0,1]\times\operatorname{Spin}(n+k)
 \xrightarrow{\ \cong\ }
 P_{\mathrm{spin}}(\mathcal L\oplus V)
 \big|_{\{p\}\times[0,1]}.
\]
$\widetilde\sigma$ 有两个选择。以
$\iota_i\circ\widetilde\sigma^{-1}$（$i=0,1$）把
$P_{\mathrm{spin}}(\mathcal L\oplus V)|_{\{p\}\times[0,1]}$
同 $P_{\mathrm{spin}}(L^{(i)},V)$ 粘合，便得到
$\mathcal L\oplus V$ 在 $K(L^{(0)},L^{(1)})_{[2]}$ 上的
spin 结构。容易看出，所得 spin 结构不依赖于 $\sigma_*$ 的提升
$\widetilde\sigma$ 之选择。

记 $\mathcal I_{\lambda_p}$ 为保持定向的
$\widetilde\lambda_p$ 平凡化 $\sigma$ 所成空间，并记
$\widetilde{\mathcal I}_{\lambda_p}$ 为上述三元组
$(\sigma,\iota_0,\iota_1)$ 所成空间。显然，
$\mathcal I_{\lambda_p}$ 同伦等价于 $SO(n)$，而
$\widetilde{\mathcal I}_{\lambda_p}$ 同伦等价于
\[
 \frac{
 \{(g_0,g_1)\in\operatorname{Spin}(n)\times\operatorname{Spin}(n)
 \mid \pi(g_0)=\pi(g_1)\}}
 {\{\pm1\}},
\]
其中 $\pi\colon\operatorname{Spin}(n)\to SO(n)$ 是典范投影，
$\{\pm1\}$ 对 $\operatorname{Spin}(n)\times\operatorname{Spin}(n)$
对角作用。因此
\[
 \widetilde{\mathcal I}_{\lambda_p}\longrightarrow
 \mathcal I_{\lambda_p}
\]
是平凡双覆盖；也就是说，
$\widetilde{\mathcal I}_{\lambda_p}$ 是
$\mathcal I_{\lambda_p}$ 的两个副本之并。

现在考虑 $\Lambda^{\mathrm{ori}}(T_pM)$ 中从 $T_pL^{(0)}$
到 $T_pL^{(1)}$ 的全部道路所成空间
$\mathcal P_p(TL^{(0)},TL^{(1)})$，并定义
\[
 \mathcal I(p)=
 \bigcup_{\lambda_p\in\mathcal P_p(TL^{(0)},TL^{(1)})}
 \mathcal I_{\lambda_p},
 \qquad
 \widetilde{\mathcal I}(p)=
 \bigcup_{\lambda_p\in\mathcal P_p(TL^{(0)},TL^{(1)})}
 \widetilde{\mathcal I}_{\lambda_p}.
\]
正如第~\ref{subsec:8-1-2} 小节所见，
$(S^1\times S^1)\times\mathbb C^n$ 中存在一个定向但非 spin 的
拉格朗日子丛。因此
\[
 \widetilde{\mathcal I}(p)\longrightarrow\mathcal I(p)
\]
是非平凡双覆盖。

对每个 $p$，选取
$\lambda_p\in\mathcal P_p(TL^{(0)},TL^{(1)})$，并考虑带帽
半无穷圆柱上的椭圆算子
\[
 \bar\partial_{\lambda_p,Z_\pm}\colon
 W_{\lambda_p}^{1,p}(Z_\pm;T_pM)
 \longrightarrow
 L^p\bigl(Z_\pm;T_pM\otimes\Lambda^{0,1}(Z_\pm)\bigr),
\]
其中
\[
\begin{aligned}
 Z_-={}&\{z\in\mathbb C\mid |z|\leq1\}
 \cup\{z\in\mathbb C\mid\operatorname{Re}z\geq0,\
       |\operatorname{Im}z|\leq1\},\\
 Z_+={}&\{z\in\mathbb C\mid\operatorname{Re}z\leq0,\
       |\operatorname{Im}z|\leq1\}
 \cup\{z\in\mathbb C\mid |z|\leq1\}.
\end{aligned}
\]
这与第~\ref{subsec:3-7-5} 小节中的定义相同；见定义
\ref{def:3-7-62} 之前。

对交点 $p\in L^{(0)}\cap L^{(1)}$ 及 $T_pM$ 中拉格朗日
子空间道路 $\lambda_p$，定义其马斯洛夫--莫尔斯指标为
\begin{equation}
 \mu(p,\lambda_p)
 =\operatorname{Index}(\bar\partial_{\lambda_p,Z_-}).
 \tag{8.1.11}\label{eq:8-1-11}
\end{equation}

对 $R>0$，如下定义空间
$Z_-\#_R(\mathbb R\times[0,1])\#_R Z_+$。

\par\medskip\noindent\begin{minipage}{\textwidth}
 \centering
 \setlength{\unitlength}{0.72mm}
 \begin{picture}(118,118)
  \linethickness{0.45pt}
  \put(30,103){\line(1,0){47}}
  \put(30,91){\line(1,0){47}}
  \qbezier(30,91)(18,91)(18,97)
  \qbezier(18,97)(18,103)(30,103)
  \multiput(40,91)(0,3){5}{\line(0,1){1.5}}
  \multiput(63,91)(0,3){5}{\line(0,1){1.5}}
  \put(40,108){\line(1,0){23}}
  \put(40,106.5){\line(0,1){3}}
  \put(63,106.5){\line(0,1){3}}
  \put(51.5,111){\makebox(0,0){$R$}}
  \put(47,87){\makebox(0,0){$Z_-$}}
  \put(36,78){\line(1,0){47}}
  \put(36,66){\line(1,0){47}}
  \multiput(48,66)(0,3){5}{\line(0,1){1.5}}
  \multiput(70,66)(0,3){5}{\line(0,1){1.5}}
  \put(48,83){\line(1,0){22}}
  \put(48,81.5){\line(0,1){3}}
  \put(70,81.5){\line(0,1){3}}
  \put(59,86){\makebox(0,0){$2R$}}
  \put(31,62){\makebox(0,0){$\mathbb R\times[0,1]$}}
  \put(43,54){\line(1,0){47}}
  \put(43,42){\line(1,0){47}}
  \qbezier(90,54)(102,54)(102,48)
  \qbezier(102,48)(102,42)(90,42)
  \multiput(65,42)(0,3){5}{\line(0,1){1.5}}
  \multiput(89,42)(0,3){5}{\line(0,1){1.5}}
  \put(65,59){\line(1,0){24}}
  \put(65,57.5){\line(0,1){3}}
  \put(89,57.5){\line(0,1){3}}
  \put(77,62){\makebox(0,0){$R$}}
  \put(35,38){\makebox(0,0){$Z_+$}}
  \put(22,36){\vector(0,-1){7}}
  \put(30,24){\line(1,0){59}}
  \put(30,12){\line(1,0){59}}
  \qbezier(30,12)(18,12)(18,18)
  \qbezier(18,18)(18,24)(30,24)
  \qbezier(89,24)(101,24)(101,18)
  \qbezier(101,18)(101,12)(89,12)
  \multiput(30,12)(0,3){5}{\line(0,1){1.5}}
  \multiput(89,12)(0,3){5}{\line(0,1){1.5}}
  \put(30,29){\line(1,0){59}}
  \put(30,27.5){\line(0,1){3}}
  \put(89,27.5){\line(0,1){3}}
  \put(59.5,32){\makebox(0,0){$4R$}}
  \put(59.5,7){\makebox(0,0)
   {$Z_-\#_R(\mathbb R\times[0,1])\#_R Z_+$}}
 \end{picture}
 \par\smallskip
 \renewcommand{\thefigure}{8.1.2}
 \refstepcounter{figure}
 \textbf{图~\thefigure}\label{fig:8-1-2}
\end{minipage}\par\medskip

考虑如下截断部分%
\translatornote{源书在第二个集合中误写 $z\in Z_-$；
由 $\operatorname{Re}z\geq-R$ 的右端帽型别及随后使用的
$Z_{+,R}$，这里恢复为 $z\in Z_+$。}：
\[
 Z_{-,R}=\{z\in Z_-\mid\operatorname{Re}z\leq R\},
 \qquad
 Z_{+,R}=\{z\in Z_+\mid\operatorname{Re}z\geq-R\}.
\]
粘合三个空间 $Z_{-,R}$、$[-R,R]\times[0,1]$、$Z_{+,R}$：
分别把 $(R,t)\in Z_{-,R}$ 与
$(-R,t)\in[-R,R]\times[0,1]$ 识别，并把
$(R,t)\in[-R,R]\times[0,1]$ 与 $(-R,t)\in Z_{+,R}$ 识别。

考虑 $Z_\pm$ 上的算子 $\bar\partial_{\lambda_p,Z_\pm}$，
以及 $\mathbb R\times[0,1]$ 上的 $D_u\bar\partial$；这里
$D_u\bar\partial$ 是柯西--黎曼方程在 $u$ 处的线性化。
利用适当的单位分解粘合这三个算子，便在
$Z_-\#_R(\mathbb R\times[0,1])\#_R Z_+$ 上得到算子
\[
 \bar\partial_{u;\lambda_{p^+},\lambda_{p^-}}.
\]
于是指标求和公式给出%
\translatornote{源书在式 (8.1.12) 左边两个帽算子中都只写
$\lambda_p$；负端帽接在 $u(-\infty,*)=p^-$，正端帽接在
$u(+\infty,*)=p^+$，且右边及式 (8.1.13) 均保留这两个端点，
故这里分别补为 $\lambda_{p^-}$、$\lambda_{p^+}$。}
\begin{equation}
\begin{aligned}
 &\operatorname{Index}(\bar\partial_{\lambda_{p^-},Z_-})
 +\operatorname{Index}(D_u\bar\partial)
 +\operatorname{Index}(\bar\partial_{\lambda_{p^+},Z_+})\\
 &\qquad={}
 \operatorname{Index}
 (\bar\partial_{u;\lambda_{p^+},\lambda_{p^-}}).
\end{aligned}
\tag{8.1.12}\label{eq:8-1-12}
\end{equation}

用适当的单位分解把 $u$ 同到 $p^\pm$ 的常值映射粘合，定义
\[
 \widehat u\colon
 Z_-\#_R(\mathbb R\times[0,1])\#_R Z_+
 \longrightarrow M.
\]
将丛平凡化为
\[
 \widehat u^*TM\cong
 \bigl(Z_-\#_R(\mathbb R\times[0,1])\#_R Z_+\bigr)
 \times\mathbb C^n.
\]
于是得到拉格朗日格拉斯曼流形中的道路
\[
 \lambda\colon
 \partial\bigl(Z_-\#_R(\mathbb R\times[0,1])\#_R Z_+\bigr)
 \longrightarrow\Lambda(\mathbb C^n),
\]
其定义为
\[
\begin{aligned}
 \lambda\bigl(e^{\pi\sqrt{-1}(1/2\pm t)}\bigr)
 &=\lambda_{p^\pm}(t),
 &e^{\pi\sqrt{-1}(1/2\pm t)}&\in\partial Z_\pm,\\
 \lambda(\tau,i)&=T_{u(\tau,i)}L^{(i)},
 &i&=0,1.
\end{aligned}
\]
容易看出
\begin{equation}
 \operatorname{Index}
 (\bar\partial_{u;\lambda_{p^+},\lambda_{p^-}})
 =\operatorname{Index}(\bar\partial_{\mathbb C^n,\lambda}),
 \tag{8.1.13}\label{eq:8-1-13}
\end{equation}
其中右边与命题~\ref{prop:8-1-4} 中相同。选取
$(\sigma,\iota_0,\iota_1)\in
\widetilde{\mathcal I}_{\lambda_p}$，则
$\mathcal L\oplus V$ 的 spin 结构以定理~\ref{thm:8-1-1}
证明中的同样方式，诱导虚向量空间
$\operatorname{Index}(\bar\partial_{\mathbb C^n,\lambda})$ 的定向。

因此，一旦选取
$\operatorname{Index}(\bar\partial_{\lambda_{p^+},Z_\pm})$ 与
$\operatorname{Index}(\bar\partial_{\lambda_{p^-},Z_\pm})$ 的定向，
以及
\[
 (\sigma^\pm,\iota_0^\pm,\iota_1^\pm)
 \in\widetilde{\mathcal I}_{\lambda_{p^\pm}},
\]
式~\eqref{eq:8-1-12}、\eqref{eq:8-1-13} 便诱导
$\operatorname{Index}(D_u\bar\partial)$ 的定向。再结合下面
注~\ref{rem:8-1-15}(1)，我们证明了如下结果。

\begin{theorem}{8.1.14}\label{thm:8-1-14}
设一对拉格朗日子流形 $(L^{(0)},L^{(1)})$ 是相对 spin 的。
那么，对任意 $p^\pm\in L^{(0)}\cap L^{(1)}$，拉格朗日交
弗洛尔上同调中连接轨道的模空间 $\mathcal M(p^+,p^-)$ 可定向。
此外，对每个 $p=p^\pm$，
$\operatorname{Index}(\bar\partial_{\lambda_p,Z_-})$ 的定向、
一个三元组
$(\sigma,\iota_0,\iota_1)\in
\widetilde{\mathcal I}_{\lambda_p}$，以及拉格朗日对子
$(L^{(0)},L^{(1)})$ 的相对 spin 结构，典范地确定
$\mathcal M(p^+,p^-)$ 的定向。
\end{theorem}

\begin{remark}{8.1.15}\label{rem:8-1-15}
\begin{enumerate}[label=(\arabic*),leftmargin=2.8em]
\item 考虑 $u_p$，其中 $u_p(\tau,t)\equiv p$；如上粘合
$\bar\partial_{\lambda_p,Z_-}$、$D_{u_p}\bar\partial$ 与
$\bar\partial_{\lambda_p,Z_+}$，得到算子
$\bar\partial_{u_p;\lambda_p,\lambda_p}$。后者的指标典范同构于
$T_pL^{(0)}$，因而带有定向；$D_{u_p}\bar\partial$ 的指标是
平凡的，也带有定向。因此，由~\eqref{eq:8-1-12}，
$\operatorname{Index}(\bar\partial_{\lambda_p,Z_-})$ 的定向决定
$\operatorname{Index}(\bar\partial_{\lambda_p,Z_+})$ 的定向。
所以，在定理~\ref{thm:8-1-14} 中，只须对每个 $p=p^\pm$
选取 $\operatorname{Index}(\bar\partial_{\lambda_p,Z_-})$ 的一个定向。%
\translatornote{源书本句在 $\operatorname{Index}$ 后又多写一层
$\bar\partial_{(\bar\partial_{\lambda_p,Z_-})}$；依紧前两句恢复为
$\operatorname{Index}(\bar\partial_{\lambda_p,Z_-})$。}

\item 相对 spin 结构所确定的定向（定理~\ref{thm:8-1-1}
与定理~\ref{thm:8-1-14}）自动满足 \textbf{[FlHo93]} 的相容性条件，
即在胶合过程中相容。关于全纯圆盘模空间定向的相容性，见下文
命题 8.3.3 与引理 8.3.5；连接轨道模空间的相容性也以引理 8.3.5
为基础得到证明，见命题 8.7.3 的证明。定理~\ref{thm:8-1-14}
中的赋予定向步骤沿用 \textbf{[FuOn99II]} 第 21 节的论证；相容定向
问题的论证另见 \textbf{[FlHo93]}。

应当指出，\textbf{[FlHo93]} 意义下的相容定向系统并不唯一。
事实上，存在会导出彼此不同构上同调的相容定向系统。在有限维情形，
若流形非单连通，可以用局部系统扭曲莫尔斯复形，得到以局部系统为
系数的上同调；一般而言，所得扭曲上同调与未扭曲上同调不同构。
局部系统的和乐会影响“连接轨道计数”中的符号。在拉格朗日子流形的
弗洛尔理论中，Cheol-Hyun Cho 找到了会导出彼此不同构弗洛尔上同调的
spin 结构 \textbf{[Cho04]}。

换言之，为了定义弗洛尔上同调，仅仅证明存在满足相容性条件的定向
并不足够；还必须指定以何种方式为各模空间赋予定向，同时使其满足
相容性条件。

\item 在 $\mathcal P_p(TL^{(0)},TL^{(1)})$ 的每个连通分支上，
$\operatorname{Index}\bar\partial_{\cdot,Z_-}$ 的行列式线丛不可定向，
但它拉回到 $\widetilde{\mathcal I}(p)$ 后可定向。为此，沿用 (1) 的
论证。选定
\[
 (\lambda_p,\sigma,\iota_0,\iota_1)
 \in\widetilde{\mathcal I}(p).
\]
把算子族 $\bar\partial_{\cdot,Z_-}$ 同
$\bar\partial_{\lambda_p,Z_+}$ 粘合，得到一族带全实边界条件的
多尔博型算子。第~\ref{subsec:8-1-2} 小节的例表明，该族指标的
行列式线丛不可定向。由于
$\bar\partial_{\lambda_p,Z_+}$ 不依赖于
$\mathcal P_p(TL^{(0)},TL^{(1)})$ 中的元素，
$\bar\partial_{\cdot,Z_-}$ 的指标丛不可定向。

容易看出，经
\[
 \mathcal I(p)\longrightarrow
 \mathcal P_p(TL^{(0)},TL^{(1)})
\]
拉回后，该丛在每个连通分支上仍不可定向，因为全实边界条件不依赖于
$\mathcal I_{\lambda_p}$ 中的元素。注意，
$\widetilde{\mathcal I}(p)\to\mathcal I(p)$ 是在
$\mathcal I(p)$ 每个连通分支上均非平凡的唯一双覆盖。因此，指标丛
进一步拉回后应当可定向，这就证明了断言。（也可以直接看出第二个
断言，因为 $\widetilde{\mathcal I}(p)$ 给出全实子丛上一族相容的
spin 结构。）

所以，“选取
$\operatorname{Index}(\bar\partial_{\lambda_{p^\pm},Z_-})$
的定向以及 $(\sigma^\pm,\iota_0^\pm,\iota_1^\pm)$”这一资料，
可以改述为：选取 $\widetilde{\mathcal I}(p)$ 上算子族
$\operatorname{Index}\bar\partial_{\cdot,Z_-}$ 的行列式线丛之定向。
后一解释更为典范，也可用于其他情形，例如博特--莫尔斯情形；
见命题 8.8.1。

\item 在 \textbf{[Oh97II]} 中，第二作者先前研究了拉格朗日交弗洛尔
同调的定向问题，其中
\[
 L^{(0)}=\phi(o_N),\qquad L^{(1)}=o_N,
 \qquad M=T^*N;
\]
更一般地，也可取 $L^{(0)}=\nu^*S$，这里 $\nu^*S$ 是
$S\subset N$ 的余法丛。他证明，在此情形下，无论紧流形 $N$ 是否
可定向，弗洛尔复形上都存在相容定向。此时关于 $w_2$ 的条件自动满足，
并且不存在边界位于 $L^{(0)}$ 或 $L^{(1)}$ 上的全纯圆盘。
\end{enumerate}
\end{remark}
